% =====================================================================
%  2 Stakeholder  (Aufgabenstellung: „Identifikation und Beschreibung der
%  anvisierten Ziel-/Benutzergruppe(n)“)
%  Tipp: Persona(s) für die Work-Life-Balance-Zielgruppe.
% =====================================================================
\section{Stakeholder}
\TODO{Kurze Hinführung}

\subsection{Stakeholderübersicht}
\begin{longtable}{@{}L{3.2cm}L{6.6cm}L{5.4cm}@{}}
  \caption{Stakeholder und ihre Interessen}\label{tab:stakeholder}\\
  \toprule
  \textbf{Stakeholder} & \textbf{Beschreibung} & \textbf{Interessen und Erwartungen} \\
  \midrule
  \endfirsthead
  \toprule
  \textbf{Stakeholder} & \textbf{Beschreibung} & \textbf{Interessen und Erwartungen} \\
  \midrule
  \endhead
  \bottomrule
  \endfoot
  \TODO{Hauptnutzergruppe} & \TODO{} & \TODO{} \\
  \TODO{weitere Gruppe} & \TODO{} & \TODO{} \\
  Tutor:in (Prüfende) & Bewertet das Portfolio und gilt im Kurs als Zielgruppe & \TODO{} \\
\end{longtable}
\quelle{Eigene Darstellung.}

\subsection{Zielgruppe und Personas}
\TODO{Zentrale Zielgruppe genauer beschreiben (Bedürfnisse, Probleme, Nutzungskontext); ggf. 1--2 Personas}
